% =============================================================================
% Example PhD progress report using the qreport class
% =============================================================================
% This is a standalone, single-file template for PhD progress (and similar)
% reports. It uses the qreport class, which provides an article-style layout
% with a compact title block and metadata for the candidate, programme, and
% reporting period.
%
% The structure and headings below are GENERIC. Adapt them to match your own
% institution's progress review guidance (for example, the COMSC PGR Review
% Guidance document) by renaming, reordering, or adding sections as needed.
% This file is only for the narrative report text, not for any administrative
% forms or approval sheets that your institution may require.
%
% If you want to record licensing information for the report itself (for
% example, a Creative Commons license), you can add a short unnumbered
% ``License'' section near the end of the document with your preferred text,
% such as a simple CC-BY or CC-BY-SA notice.
%
% Build with lualatex, xelatex, or pdflatex. No other project files are needed
% apart from an optional .bib file if you use \bib{...} at the end.
%
% SPDX-FileCopyrightText: 2026 Frank Langbein <frank@langbein.org>
% SPDX-License-Identifier: LPPL-1.3c
%

\documentclass[sans,oneside,numeric,raggedright]{qreport}

% -----------------------------------------------------------------------------
% Report metadata
% -----------------------------------------------------------------------------

\title{PhD Progress Report}

% Candidate and programme
\candidate{Your Name}
\studentid{1234567}
\programme{PhD in Computer Science}
\institution{Your University}
\school{School of Computer Science}

% Supervisors and report details
\supervisors{Supervisor One, Supervisor Two}
\reporttype{Annual Progress Review}
\reportperiod{1 October 2025 -- 30 September 2026}

% Date (defaults to \today if omitted)
\date{\today}

% Optional: keywords appear at the end of the abstract environment
% (if you use \maketitlewithabstract)
\keywords{progress report; PhD; research degree}

% Optional: useful packages for reports (tables, math, etc.)
\usepackage{booktabs}
\usepackage{mathtools}
\usepackage{siunitx}
\usepackage{enumitem}

\begin{document}

% Compact, article-style title at top of first page
\maketitle

% Optionally, you can replace the above with:
% \maketitlewithabstract{%
%   Provide a short summary of the reporting period: key achievements,
%   main results, and overall status. Keep this to one paragraph if your
%   institution expects a brief overview at the start of the report.
% }

% =============================================================================
% Suggested section structure
% =============================================================================
% The sections below reflect a typical PhD progress report. Adapt headings,
% order, and level of detail to match your local requirements and review
% forms. You can merge or split sections (e.g. combining Outputs with
% Research Activities) as appropriate.

\section{Summary of progress}

\newthought{This section provides a concise overview of your year.}
Summarise the reporting period in a few paragraphs. Highlight the main
achievements since the last report, your current overall status (on track,
ahead, behind schedule), and any major changes in direction or scope.

\section{Research activities and results}

Describe the research you have carried out during the period. Depending on
your project, you might organise this by research question, work package,
experiment, or theoretical component.

\subsection{Work completed}

Outline the main tasks completed (e.g.\ literature study, theoretical
development, software implementation, experimental design, data collection,
analysis). Be specific enough that reviewers can see clear progress.

\subsection{Results and findings}

Summarise the key findings so far. This can include preliminary results,
negative results, or partial analyses. If appropriate, refer to figures,
tables, or appendices that document results more fully.

\section{Outputs and dissemination}

List publications, submissions, preprints, talks, posters, software, and
datasets arising from your work during this period (and, if helpful, earlier
work that is directly relevant to the thesis). You can use an itemised list
or a table.

\begin{itemize}[leftmargin=2em]
  \item Journal article or conference paper (published, accepted, submitted,
        or in preparation), with a short note on status.
  \item Talks, posters, and demos given at conferences, seminars, or internal
        events.
  \item Software, code, or data releases (e.g.\ Git repositories, DOIs).
\end{itemize}

If you maintain a bibliography file, you can also cite outputs here using
standard \texttt{natbib} commands (e.g.\ \verb|\citep|, \verb|\citet|).

\section{Training and professional development}

Describe any formal or informal training undertaken during the period.
Typical examples include:

\begin{itemize}[leftmargin=2em]
  \item Research methods courses and technical workshops.
  \item Generic skills training (writing, presenting, teaching, project
        management, statistics, etc.).
  \item Teaching, demonstrating, or supervision activities.
  \item Outreach, public engagement, or professional service.
\end{itemize}

Comment briefly on how these activities support your research and your
longer-term development.

\section{Plan for the next period}

Provide a concrete plan for the next reporting period. Align this plan
with your thesis timeline or Gantt chart if you have one.

\begin{itemize}[leftmargin=2em]
  \item State your main objectives and milestones for the next period.
  \item Indicate any dependencies or sequencing (e.g.\ data collection
        before analysis, method development before experiments).
  \item Note where work will continue existing strands versus starting new
        directions.
\end{itemize}

You can optionally include a table or bullet list mapping planned tasks to
time periods (months or semesters).

\section{Risks, challenges, and mitigation}

Discuss any risks or challenges affecting your progress. These might include
technical difficulties, resource constraints, access to data or equipment,
changes in supervision, or personal circumstances (describe only at the level
of detail you are comfortable sharing in an academic report).

For each significant risk or challenge, outline:

\begin{itemize}[leftmargin=2em]
  \item Its impact on your project and timeline.
  \item Steps already taken to address it.
  \item Planned mitigation or contingency measures.
\end{itemize}

\section{Additional information}

Include here any information that does not fit naturally in the earlier
sections but is requested by your institution's progress review forms (for
example, ethical approvals, data management plans, or specific questions
about supervision or resources).

You can also use this section to record agreed changes to your project
scope or thesis structure, if these are part of the review.

% =============================================================================
% Optional bibliography
% =============================================================================
% If your institution expects references in the progress report (for example,
% when you cite key background work or your own published outputs), include a
% .bib file and uncomment the lines below. Replace ``references'' with the
% name of your .bib file (without the .bib extension).
%
% \bibliographystyle{plainnat} is selected automatically by \bib{...} for
% numeric citations, or agsm for harvard-style citations.

% \clearpage
% \bib{references}

% =============================================================================
% Optional license information
% =============================================================================
% If you need to state a license for the report content (for example, that it
% is released under a Creative Commons license), you can add a short unnumbered
% section such as:
%
% \section*{License}
% This report is licensed under the Creative Commons Attribution-ShareAlike 4.0
% International License (CC BY-SA 4.0). You are free to share and adapt the
% material under the terms of that license.
%
% Adjust the wording and license choice to match your needs or your
% institution's recommendations.

\end{document}

